\subtitle{\bf Aula 2 - Análise Léxica}
\begin{frame}[t,plain]
\titlepage
\end{frame}

\begin{frame}{Objetivos}
\protect\hypertarget{objetivos-1}{}
\begin{itemize}
\item
  Apresentar a importância da análise léxica.
\item
  Apresentar a implementação de um analisador léxico ad-hoc.
\end{itemize}
\end{frame}



\begin{frame}{O Papel da Análise Léxica}
\protect\hypertarget{o-papel-da-anuxe1lise-luxe9xica}{}
\begin{itemize}

\item
  Primeira fase do compilador
\item
  Recebe o código-fonte como \textbf{sequência de caracteres}
\item
  Agrupa os caracteres em \textbf{tokens} significativos
\end{itemize}
\end{frame}

\begin{frame}{O Papel da Análise Léxica}
\protect\hypertarget{o-papel-da-anuxe1lise-luxe9xica-1}{}
\begin{itemize}

\item
  Remove espaços em branco e comentários
\item
  Detecta erros léxicos precocemente (caracteres inválidos, etc.)
\item
  Funciona como a ``porta de entrada'' que organiza o texto bruto
\end{itemize}
\end{frame}

\begin{frame}{O que é um Token?}
\protect\hypertarget{o-que-uxe9-um-token}{}
\begin{itemize}

\item
  Um \textbf{token} é um componente indivisível da sintaxe de uma
  linguagem
\item
  O conjunto de tokens pode ser entendido como o \textbf{alfabeto} da
  linguagem
\item
  Exemplos de tokens: palavras reservadas, identificadores, números,
  operadores
\end{itemize}
\end{frame}

\begin{frame}{O que é um Token?}
\protect\hypertarget{o-que-uxe9-um-token-1}{}
\begin{itemize}

\item
  O analisador léxico produz uma \textbf{lista de tokens} para o
  analisador sintático
\item
  Cada token carrega sua \textbf{categoria} (lexema) e posição no
  código-fonte
\end{itemize}
\end{frame}

\begin{frame}[fragile]{Estrutura de um Token}
\protect\hypertarget{estrutura-de-um-token}{}
\begin{lstlisting}[language=Haskell]
type Line   = Int
type Column = Int

data Token = Token
  { pos    :: (Line, Column)  -- posição no código-fonte
  , lexeme :: Lexeme          -- categoria do token
  }
\end{lstlisting}

\begin{itemize}

\item
  A posição é essencial para \textbf{mensagens de erro} precisas
\item
  O lexema identifica a categoria do token
\end{itemize}
\end{frame}

\hypertarget{a-linguagem-exp}{%
\section{A Linguagem Exp}\label{a-linguagem-exp}}

\begin{frame}{Gramática da Linguagem Exp}
\protect\hypertarget{gramuxe1tica-da-linguagem-exp}{}
\[E \to n \mid E + E \mid E * E \mid (E)\]

\begin{itemize}

\item
  \(n\): constante numérica inteira
\item
  \(+\): operador de adição
\item
  \(*\): operador de multiplicação
\item
  \((\) e \()\): parênteses para agrupamento
\item
  Linguagem simples: expressões aritméticas básicas
\end{itemize}
\end{frame}

\begin{frame}[fragile]{Tokens da Linguagem Exp}
\protect\hypertarget{tokens-da-linguagem-exp}{}
\begin{table}
\begin{tabular}[]{@{}lll@{}}
\toprule
Token & Construtor & Descrição \\
\midrule

\texttt{n} & \texttt{TNumber Int} &
número inteiro (carrega o valor) \\
\texttt{+} & \texttt{TPlus} & operador
de adição \\
\texttt{*} & \texttt{TTimes} &
operador de multiplicação \\
\bottomrule
\end{tabular}
\end{table}
\end{frame}

\begin{frame}[fragile]{Tokens da Linguagem Exp}
\protect\hypertarget{tokens-da-linguagem-exp-1}{}
\begin{table}
\begin{tabular}[]{@{}lll@{}}
\toprule
Token & Construtor & Descrição \\
\midrule

\texttt{(} & \texttt{TLParen} & abre
parênteses \\
\texttt{)} & \texttt{TRParen} & fecha
parênteses \\
EOF & \texttt{TEOF} & marcador de fim de entrada \\
\bottomrule
\end{tabular}
\end{table}
\end{frame}

\begin{frame}[fragile]{Definição do Tipo Lexeme}
\protect\hypertarget{definiuxe7uxe3o-do-tipo-lexeme}{}
\begin{lstlisting}[language=Haskell]
data Lexeme
  = TNumber Int   -- número com seu valor inteiro
  | TLParen       -- (
  | TRParen       -- )
  | TPlus         -- +
  | TTimes        -- *
  | TEOF          -- fim de entrada
\end{lstlisting}
\end{frame}

\begin{frame}[fragile]{Definição do Tipo Lexeme}
\protect\hypertarget{definiuxe7uxe3o-do-tipo-lexeme-1}{}
\begin{itemize}

\item
  Cada construtor corresponde a uma categoria de token
\item
  \texttt{TNumber} carrega o valor inteiro do número
\item
  Os demais construtores apenas identificam o símbolo
\end{itemize}
\end{frame}

\hypertarget{analisadores-luxe9xicos-ad-hoc}{%
\section{Analisadores Léxicos
Ad-Hoc}\label{analisadores-luxe9xicos-ad-hoc}}

\begin{frame}{Analisadores Ad-Hoc}
\protect\hypertarget{analisadores-ad-hoc}{}
\begin{itemize}

\item
  Implementação \textbf{manual} e específica para uma linguagem
\item
  Processa caracteres sequencialmente com casamento de padrões
\end{itemize}
\end{frame}

\begin{frame}{Analisadores Ad-Hoc}
\protect\hypertarget{analisadores-ad-hoc-1}{}
\begin{itemize}

\item
  Não usa geradores automáticos como o Alex
\item
  Adequado para linguagens \textbf{pequenas e simples}
\item
  Oferece controle total sobre as mensagens de erro
\end{itemize}
\end{frame}

\begin{frame}[fragile]{Tipo do Analisador}
\protect\hypertarget{tipo-do-analisador}{}
\begin{lstlisting}[language=Haskell]
-- Versão simples (sem tratamento de erros):
lexer :: String -> [Token]

-- Versão com tratamento de erros:
lexer :: String -> Either String [Token]
\end{lstlisting}

\begin{itemize}

\item
  \texttt{Left msg} indica um \textbf{erro léxico} com
  mensagem \texttt{msg}
\item
  \texttt{Right tokens} indica \textbf{sucesso} com a
  lista de tokens produzida
\end{itemize}
\end{frame}

\begin{frame}[fragile]{O Tipo Either}
\protect\hypertarget{o-tipo-either}{}
\begin{lstlisting}[language=Haskell]
data Either a b = Left a | Right b
\end{lstlisting}

\begin{itemize}

\item
  \texttt{Left a} representa falha/erro com valor do
  tipo \texttt{a}
\item
  \texttt{Right b} representa sucesso com valor do tipo
  \texttt{b}
\end{itemize}
\end{frame}

\begin{frame}[fragile]{Estado do Analisador}
\protect\hypertarget{estado-do-analisador}{}
\begin{lstlisting}[language=Haskell]
type State = (Line, Column, String, [Token])
--            ^      ^       ^        ^
--           linha  coluna  dígitos  tokens
--                          acum.    produzidos
\end{lstlisting}

\begin{itemize}

\item
  \texttt{Line}, \texttt{Column}:
  posição atual no arquivo de entrada
\item
  \texttt{String}: dígitos consecutivos acumulados
  (para montar números)
\item
  \texttt{[Token]}: tokens reconhecidos até o momento
\end{itemize}
\end{frame}

\begin{frame}[fragile]{Função de Transição}
\protect\hypertarget{funuxe7uxe3o-de-transiuxe7uxe3o}{}
\begin{lstlisting}[language=Haskell]
transition :: State -> Char -> Either String State
transition state@(l, col, t, ts) c
  | c == '\n'  = mkDigits state c   -- nova linha
  | isSpace c  = mkDigits state c   -- espaço
  | c == '+'   = Right (l, col+1, "", mkToken state TPlus    : ts)
  | c == '*'   = Right (l, col+1, "", mkToken state TTimes   : ts)
  | c == '('   = Right (l, col+1, "", mkToken state TLParen  : ts)
  | c == ')'   = Right (l, col+1, "", mkToken state TRParen  : ts)
  | isDigit c  = Right (l, col+1, c:t, ts)  -- acumula dígito
  | otherwise  = unexpectedCharError l col c
\end{lstlisting}
\end{frame}

\begin{frame}[fragile]{Tratamento de Dígitos}
\protect\hypertarget{tratamento-de-duxedgitos}{}
\begin{lstlisting}[language=Haskell]
mkDigits :: State -> Char -> Either String State
mkDigits state@(l, col, s, ts) c
  | null s      = -- nenhum dígito acumulado: só atualiza posição
      let l'   = if c == '\n' then l + 1 else l
          col' = if isSpace c then col + 1 else col
      in Right (l', col', s, ts)
  | all isDigit s = -- dígitos válidos: cria token TNumber
      let t    = Token (l,col) (TNumber (read $ reverse s))
          l'   = if c == '\n' then l + 1 else l
          col' = if c /= '\n' && isSpace c then col + 1 else col
      in Right (l', col', "", t : ts)
  | otherwise   = unexpectedCharError l col c
\end{lstlisting}
\end{frame}

\begin{frame}[fragile]{A Função \texttt{lexer}
Completa}
\protect\hypertarget{a-funuxe7uxe3o-lexer-completa}{}
\begin{lstlisting}[language=Haskell]
lexer :: String -> Either String [Token]
lexer = finish . foldl step base
  where
    base          = Right (1, 1, "", [])
    finish        = either Left (Right . extract)
    step (Left e) _ = Left e
    step (Right s) c = transition s c

    extract (l, col, s, ts)
      | null s    = reverse ts
      | otherwise = let n = read (reverse s)
                        t = Token (l, col - length s) (TNumber n)
                    in reverse (t : ts)
\end{lstlisting}
\end{frame}

\begin{frame}[fragile]{Como \texttt{foldl} Funciona}
\protect\hypertarget{como-foldl-funciona}{}
\begin{lstlisting}[language=Haskell]
foldl :: (b -> a -> b) -> b -> [a] -> b
foldl _ v []     = v
foldl f v (x:xs) = foldl f (f v x) xs
\end{lstlisting}

Exemplo: \texttt{sum [1,2,3]} com
\texttt{foldl (+) 0}:

\begin{lstlisting}
foldl (+) 0 [1,2,3]
= foldl (+) (0+1) [2,3]
= foldl (+) (0+1+2) [3]
= foldl (+) (0+1+2+3) []
= 6
\end{lstlisting}

\begin{itemize}

\item
  \texttt{foldl} percorre a lista \textbf{da esquerda
  para a direita}
\item
  Acumula um resultado a cada passo
\end{itemize}
\end{frame}

\begin{frame}[fragile]{Exemplo de Execução}
\protect\hypertarget{exemplo-de-execuuxe7uxe3o}{}
Entrada: \texttt{"2 + 3"}

\begin{table}
\begin{tabular}[]{@{}ll@{}}
\toprule
Char & Estado (linha, col, dígitos, tokens) \\
\midrule

\texttt{2} &
\texttt{(1, 2, "2", [])} \\
\texttt{\_} &
\texttt{(1, 3, "", [TNumber 2])} \\
\texttt{+} &
\texttt{(1, 4, "", [TPlus, TNumber 2])} \\
\bottomrule
\end{tabular}
\end{table}
\end{frame}

\begin{frame}[fragile]{Exemplo de Execução}
\protect\hypertarget{exemplo-de-execuuxe7uxe3o-1}{}
Entrada: \texttt{"2 + 3"}

\begin{table}
\begin{tabular}[]{@{}ll@{}}
\toprule
Char & Estado (linha, col, dígitos, tokens) \\
\midrule

\texttt{\_} &
\texttt{(1, 5, "", [TPlus, TNumber 2])} \\
\texttt{3} &
\texttt{(1, 6, "3", [TPlus, TNumber 2])} \\
EOF & extrai ->
\texttt{[TNumber 2, TPlus, TNumber 3]} \\
\bottomrule
\end{tabular}
\end{table}
\end{frame}

\hypertarget{vantagens-e-desvantagens}{%
\section{Vantagens e Desvantagens}\label{vantagens-e-desvantagens}}

\begin{frame}{Vantagens dos Analisadores Ad-Hoc}
\protect\hypertarget{vantagens-dos-analisadores-ad-hoc}{}
\begin{itemize}

\item
  \textbf{Simples de implementar} para linguagens pequenas
\item
  \textbf{Controle total} sobre mensagens de erro
\item
  Sem dependências externas (geradores de código)
\end{itemize}
\end{frame}

\begin{frame}{Vantagens dos Analisadores Ad-Hoc}
\protect\hypertarget{vantagens-dos-analisadores-ad-hoc-1}{}
\begin{itemize}

\item
  Fácil de entender e depurar para casos simples
\item
  Pode ser altamente otimizado para o caso específico
\end{itemize}
\end{frame}

\begin{frame}[fragile]{Desvantagens dos Analisadores Ad-Hoc}
\protect\hypertarget{desvantagens-dos-analisadores-ad-hoc}{}
\begin{itemize}

\item
  \textbf{Difíceis de manter e extender}:

  \begin{itemize}

  \item
    Como adicionar identificadores?
  \item
    Como suportar comentários de linha (\texttt{//})? E
    comentários em bloco (\texttt{/* */})?
  \end{itemize}
\end{itemize}
\end{frame}

\begin{frame}{Desvantagens dos Analisadores Ad-Hoc}
\protect\hypertarget{desvantagens-dos-analisadores-ad-hoc-1}{}
\begin{itemize}

\item
  \textbf{Ineficientes para estruturas léxicas complexas}

  \begin{itemize}

  \item
    Cada nova funcionalidade exige modificar toda a estrutura do estado
  \item
    Não há garantia de correção: é fácil esquecer um caso
  \end{itemize}
\end{itemize}
\end{frame}

\begin{frame}[fragile]{Dificuldades}
\protect\hypertarget{dificuldades}{}
\begin{lstlisting}
if (x > 0) ift = x + 1;
\end{lstlisting}

\begin{itemize}

\item
  Como distinguir \texttt{if} (palavra reservada) de
  \texttt{ift} (identificador)?
\item
  Como tratar palavras reservadas de forma eficiente?
\end{itemize}
\end{frame}

\begin{frame}{Dificuldades}
\protect\hypertarget{dificuldades-1}{}
\begin{itemize}

\item
  Para linguagens reais: dezenas de tokens, comentários aninhados, etc.
\item
  Solução: \textbf{expressões regulares} + \textbf{autômatos finitos}
\end{itemize}
\end{frame}

\hypertarget{conclusuxe3o}{%
\section{Conclusão}\label{conclusuxe3o}}

\begin{frame}[fragile]{Conclusão}
\protect\hypertarget{conclusuxe3o-1}{}
\begin{itemize}

\item
  Análise léxica transforma \textbf{texto bruto} em uma \textbf{lista de
  tokens}
\item
  Cada token carrega categoria (lexema) e posição no fonte
\item
  O tipo \texttt{Either} de Haskell permite tratamento
  elegante de erros
\end{itemize}
\end{frame}

\begin{frame}{Conclusão}
\protect\hypertarget{conclusuxe3o-2}{}
\begin{itemize}

\item
  Analisadores ad-hoc funcionam para linguagens simples
\item
  As limitações motivam o uso de \textbf{expressões regulares e
  autômatos}
\end{itemize}
\end{frame}
